\chapter{Родительский вывод носителя спинорного накрытия и калибровочный критерий остановки}

Проверка функции склейки доказала, что пространственная граница Каллиаса и Real-Toeplitz
границы несут один класс $+15/-15$. Остался не топологический, а
представительный вопрос: на каком физическом модуле действует масса
$n\cdot\sigma$ и выводится ли этот модуль уже принятой архитектурой.

\section{Почему пространственной двойки недостаточно}

Пусть пространственный оператор Дирака действует на
$\mathbb C^2_{\rm spin}$ матрицами $\sigma_i$. В принятой
Dirac--Schrödinger форме потенциал Каллиаса должен действовать на
независимом скручивающем множителе и коммутировать с пространственным
Clifford-действием. Именно это превращает коммутатор с оператором Дирака
в нуль-порядковый член. Такая гипотеза используется в индексных теоремах
Ангела; см. \path{Commun.Math.Phys.128(1990)77-97} и современную
формулировку \path{arXiv:2407.02993}.

Если попытаться использовать одну и ту же двойку,
\begin{equation}
 c(dx^i)=\sigma_i,
 \qquad
 \Phi=n_a\sigma_a,
\end{equation}
то
\begin{equation}
 [\sigma_i,n_a\sigma_a]\ne0
\end{equation}
для общих $i,n$. Поэтому $[D,\Phi]$ содержит первый порядок. Машинный
аудит дал максимальную норму такого матричного коммутатора $2\sqrt2$.

Минимальная корректная факторизация имеет вид
\begin{equation}
 \mathcal S_{\rm sp}\otimes\mathcal W,
 \qquad
 c(dx^i)=\sigma_i\otimes I_2,
 \qquad
 \Phi=I_2\otimes n_a\tau_a,
 \label{eq:v6-independent-callias-doublet}
\end{equation}
где все перекрёстные коммутаторы равны нулю. Таким образом, необходимы
две разные комплексные двойки: пространственный спин и внутренний twist.

\section{Топологическая двойка действительно существует}

Проект уже вывел сопряжённые хопфовы линии $L/L^*$. Их сумма имеет
функцию склейки
\begin{equation}
 g_{L\oplus L^*}(z)=
 \begin{pmatrix}z&0\\0&z^{-1}\end{pmatrix}
 \in SU(2).
 \label{eq:v6-hopf-pair-su2-transition}
\end{equation}
Полный первый класс Черна равен нулю. Поскольку
$\pi_1(SU(2))=0$, эта комплексная сумма ранга два топологически тривиальна
и может быть представлена как обычная $\mathbb C^2$ с меняющимися
собственными линиями $L$ и $L^*$.

Следовательно, препятствия на уровне векторного расслоения нет. Более
того, этот факт объясняет, почему локальная матрица $n\cdot\sigma$ и
глобальное разложение на $L/L^*$ являются двумя описаниями одного объекта.

\begin{proposition}[Топологический носитель без физической
эквивариантности]
Сумма $L\oplus L^*$ даёт необходимое скручивание ранга два как обычное
комплексное расслоение. Но для физического продолжения через ядро его две
компоненты должны допускать комплексно-линейное смешивание, совместимое со
всеми внутренними калибровочными действиями и градуировкой.
\end{proposition}

\section{Почему существующая KO6-пара не является этой двойкой}

Естественно попытаться использовать уже имеющиеся сопряжённые половины
$E/E^*$ или пакет частицы и сопряжённой частицы. На коэффициентном ранге пятнадцать
их градуировка имеет вид
\begin{equation}
 \Gamma=\tau_3\otimes I_{15}.
\end{equation}
Два недиагональных генератора предполагаемой внутренней $SU(2)$ равны
$\tau_1\otimes I_{15}$ и $\tau_2\otimes I_{15}$. Аудит даёт
\begin{equation}
 \|[\Gamma,\tau_1\otimes I_{15}]\|_F
 =\|[\Gamma,\tau_2\otimes I_{15}]\|_F
 =2\sqrt{30}=10.954451\ldots,
\end{equation}
тогда как только $\tau_3$ сохраняет градуировку.

То же препятствие видно по зарядам. Для любой ненулевой абелевой метки
сопряжённая пара несёт действие, пропорциональное
\begin{equation}
 Y=q\,\tau_3\otimes I_{15}.
\end{equation}
Недиагональные $\tau_1,\tau_2$ не коммутируют с $Y$ и смешивают
представление с его комплексно-сопряжённым. Антилинейный оператор $J$
обменивает эти половины, но не является отсутствующим комплексно-линейным
интертвинером.

Это воспроизводит, теперь непосредственно в задаче Каллиаса, ранний
представительный запрет семейного $SU(2)$ на $H_{15}/M_{35}$. Матричная
вещественность проектора $q_0$ также не спасает ситуацию: вещественный
проектор и вещественный тип полного калибровочного представления --- не
одно и то же.

\begin{nogocriterion}[Запрет смешения частицы и сопряжённой частицы]
Нельзя читать KO6-удвоение как внутренний дублет спинорного накрытия, если
недиагональные генераторы не коммутируют с калибровочной алгеброй или
градуировкой. Антилинейный Real-обмен не заменяет комплексно-линейное
$SU(2)$-действие.
\end{nogocriterion}

\section{Аудит остальных существующих кандидатов}

\paragraph{Слабый дублет.}
Центр слабого $SU(2)_L$ действует знаком минус только на восьми левых
состояниях $H_{15}$ и знаком плюс на семи правых. Поэтому он не даёт
единую двойку для коэффициентного ранга пятнадцать и изменение его роли
нарушило бы уже зафиксированные представления Стандартной модели.

\paragraph{Ориентационная двойка локального переноса.}
Она была введена как пара встречных направлений, непрерывный предел
которой даёт дираковскую кинематику. Если считать её пространственным
spin-фактором, повторное использование запрещено коммутатором выше. Если
объявить её независимым внутренним множителем, это новая внешняя двойка,
не содержащаяся в $H_{15}/M_{35}$ и не нормированная следом $M_{35}$.

\paragraph{Clifford-стабилизация.}
Добавление матричного Clifford-модуля всегда возможно для представления
$K$-класса, но стабильный аналитический представитель не является
автоматическим выводом новых физических состояний.

\paragraph{Независимая копия с теми же зарядами.}
Модуль
\begin{equation}
 \mathbb C^2_{\rm twist}\otimes H_{15}
\end{equation}
решает алгебраическую задачу: внутренние матрицы Паули коммутируют с
одинаковым калибровочным действием на двух копиях. Но число частичных
комплексных состояний меняется $15\to30$, а Real-завершение ---
$30\to60$. Это новый модуль, новый след и новый аномальный бюджет.

Стандартный монопольный механизм Джеккива--Ребби действительно использует
фермион в нетривиальном внутреннем дублете; индексная теорема Каллиаса
классифицирует его нулевые моды, но не создаёт этот дублет из ничего; см.
\path{Phys.Rev.D13(1976)3398} и
\path{Commun.Math.Phys.62(1978)213-234}.

\section{Вердикт и последняя проектная лазейка}

Получено точное разделение:
\begin{equation}
 \boxed{
 \begin{gathered}
 L\oplus L^*:\ \text{топологический носитель ранга два существует},\\
 H_{15}\oplus JH_{15}:\ \text{не является допустимым внутренним дублетом},\\
 \mathbb C^2_{\rm twist}\otimes H_{15}:\ \text{работает, но является новым модулем}.
 \end{gathered}}
\end{equation}

Таким образом, классовый мост предыдущей проверки сохраняется полностью,
но пятнадцать физических нулевых мод из текущего конечного родителя ещё
не следуют.

Остаётся одна проектная лазейка, которую нельзя отбросить без проверки.
Том VI уже построил тензорный квадрат, SWAP и носители внешних степеней из двух
копий исходного объекта. Следующая проверка должна проверить, появляется ли
там каноническая двухмерная кратность, на которой Pauli-действие
коммутирует с физической алгеброй, не удваивая фундаментальные виды
частиц. Если такое пространство кратности отсутствует, фермионная ветвь
фермионная ветвь текущей архитектуры закрывается.

\begin{protocol}
\begin{enumerate}
 \item одна $\mathbb C^2$ для пространственного спина и массы Каллиаса: \textbf{не проходит};
 \item минимальный продукт spin и twist: \textbf{размерность четыре};
 \item топологический пакет $L\oplus L^*$: \textbf{существует};
 \item его полный $c_1$: \textbf{нуль};
 \item существующая KO6-пара как комплексно-линейный $SU(2)$-дублет:
 \textbf{не проходит};
 \item слабый дублет на всех пятнадцати состояниях: \textbf{нет};
 \item ориентационная двойка как уже независимый внутренний фактор:
 \textbf{не доказана};
 \item новая равнозарядная двойка: \textbf{алгебраически работает};
 \item её цена: \textbf{новые состояния, след и аномальный аудит};
 \item физические пятнадцать нулевых мод: \textbf{не выведены};
 \item следующая проверка: \textbf{двухкопийная кратность спинорного накрытия}.
\end{enumerate}
\end{protocol}

Машинный сертификат сохранён в
\path{s2t_v6_spin_cover_carrier_parent_derivation_gate_results.json}.